
\subsubsection{Image Classification}
Describe data prep and model architecture


In our framework, we employ multiple approaches to generate counterfactual paths, moving beyond simple shortest-distance methods to incorporate more sophisticated path-finding strategies that ensure both feasibility and actionability. \todo[inline]{provide figure showing comparison}

\paragraph{Feasible Path Generation}
Drawing from the work of Sokol et al. on navigating explanatory multiverse through counterfactual path geometry~\cite{sokol2023navigating}, we implement feasible path generation. This approach aims to address two critical limitations of traditional counterfactual methods: ensuring the counterfactual instance lies within high-density regions of the data distribution, and guaranteeing the existence of a plausible transformation path from the original instance.

The feasible path generation process utilizes density-weighted graph construction, where edge weights incorporate both distance metrics and local data density information. This approach ensures that generated paths traverse high-density regions of the feature space, making the resulting counterfactuals more representative of the underlying data distribution~\cite{poyiadzi2020face}. In this work we aim to evaluate how useful these counterfactuals are compared to counterfactuals generated by the shortest path baselines.

\paragraph{Shortest Path Baselines}

To evaluate the effectiveness of feasible path strategies, we implement two baseline approaches based on shortest-distance metrics: raw Euclidean distance and PCA-transformed distance. These baselines provide comparison points for understanding when sophisticated path-finding offers advantages over simpler geometric approaches.


\paragraph{Feature Representation}
Images are provided as flattened grayscale inputs derived from the original ($28 \times 28$) MNIST digits, with pixel intensities normalised to the range ($[0, 1]$).

\paragraph{Model Architecture}
For image classification, we employed a \textbf{Convolutional Neural Network (CNN)} implemented in \texttt{JAX/Flax}. The model, denoted \texttt{MNISTConvClassifierGemini}, extracts spatial features from grayscale images and maps them to class predictions. Its architecture consists of the following components:
\begin{enumerate}
\item \textbf{Convolutional feature extractor:} Two convolutional layers with 32 and 64 filters (kernel size $3 \times 3$) are each followed by ReLU activation and max-pooling ($2 \times 2$).
\item \textbf{Feature embedding:} The flattened feature maps are projected into a 512-dimensional dense representation.
\item \textbf{Dropout regularization:} A dropout layer (rate = 0.4) is applied to the embedding during training to mitigate overfitting.
\item \textbf{Classification head:} A final dense layer maps the embedding to ten output logits corresponding to the MNIST digit classes.
\end{enumerate}
Convolutional layers use LeCun-normal initialization, while dense layers use Glorot-normal initialization. The model outputs both class logits and the intermediate feature embeddings for use in downstream analyses. All hyperparameters are optimised on the validation set for experiments with cross-entropy.